\documentclass[11pt]{article}
\usepackage[T1]{fontenc}
\usepackage{lmodern,microtype,amsmath,amssymb,array,longtable}
\usepackage[margin=1in]{geometry}
\usepackage{xurl,hyperref}
\hypersetup{colorlinks=true,linkcolor=blue,urlcolor=blue,
pdftitle={Syntactic realization: proof and review record -- Version 19}}
\newcommand{\BT}{\mathcal B_T}
\title{Syntactic realization: proof and review record}
\author{}
\date{Version 19\\1 November 2026}
\begin{document}
\maketitle

Theorem 5.9 of \emph{Local connectedness in recursion categories and
essential incompleteness}, version 19, establishes the first-order
realization in
\[
 \BT=\operatorname{Tot}(\operatorname{Split}(S_T)).
\]
Here $T$ is a consistent extension of PA. The variable and coefficient
objects are arbitrary objects of $\BT$; the number of sorts and the
arity bound are finite and fixed externally. For every axiom subobject
of the sentence object, the theorem constructs the least sentence
theory and represents it by a restriction idempotent in the original
category $S_T$. The realization is in Section 5. The mathematical
content of the abstract syntax appendix is unchanged.

\section*{Outline and dependencies}
\begin{longtable}{@{}>{\raggedright\arraybackslash}p{.18\textwidth}>{\raggedright\arraybackslash}p{.34\textwidth}>{\raggedright\arraybackslash}p{.42\textwidth}@{}}
\textbf{Location}&\textbf{Result}&\textbf{Function in the proof}\\\hline
5.1; Lemmas 5.2--5.3&Functional search and the total splitting&
Identifies $\BT$ with $\Sigma_1$ predicates and $T$-provably total
functional relations. Proves finite limits, regularity, stable disjoint
coproducts, and partial inverses for every monomorphism.\\[5pt]
5.2; Lemma 5.4&Finite witness collection&
Keeps membership and evaluation relations $\Sigma_1$ for internally
finite lists of unbounded length.\\[5pt]
5.2; Propositions 5.5--5.6&Initial syntax and lists&
Proves initiality in every slice, including algebras on arbitrary
projections to the parameter object.\\[5pt]
5.3; Proposition 5.7&Free variables and substitution&
Identifies numerical algorithms with the appendix's folds and
admissibility pullback by their PA-provable equations.\\[5pt]
5.4; Lemma 5.8&Initial derivation algebra&
Proves the universal property for every bounded-arity inference
presentation in $\BT/F$.\\[5pt]
5.4; Theorem 5.9&Generated sentence theories&
Applies Appendix B's image theorem and image--range lemma; returns
the representing range idempotent to $S_T$.\\
\end{longtable}

The decisive induction is internal to $T$. An arbitrary algebra has
operations whose relative totality and equations are theorems of $T$.
A $\Sigma_1$ evaluation relation records values and computation traces
at the nodes of a coded finite tree. Induction proves existence,
functionality, preservation of the projection to the parameter, and
the constructor equations. A second induction proves uniqueness as
uniform $T$-provable graph equality. Membership witnesses are existential
in the carrier predicate and are not additional syntax labels.

\section*{Hostile-referee review}
\subsection*{(A) Kill-shots / hard blockers}
No mathematical blocker was found in the three contributions or their
integrated realization proof. This is an author-side source review,
not a formal verification.

\subsection*{(B) Major weaknesses}
No major gap was retained after checking the following inferences.
First, monicity is tested on the internal kernel pair, yielding a
$T$-proof of injectivity; standard points are not used. Second,
preimage search splits every provably surjective morphism, giving
kernel-pair coequalizers and pullback-stable regular epimorphisms.
Third, the initiality proofs quantify over arbitrary slice algebras.
Fourth, the derivation constructor domain is the actual pullback
specified by the inference endofunctor.

\subsection*{(C) Minor issues}
The integrated text was corrected to distinguish arithmetic carrier
predicates from their split objects, numerical functions from categorical
morphisms, and the external arity bound from the derivation-code variable.
The sentence predicate is explicitly defined on the same formula-code
coordinate used by the conclusion graph.

\subsection*{(D) Cross-section consistency failures}
None found in the checked dependencies. The final theorem uses sentence
axioms, as Appendix B requires. It constructs the sentence theory by
pulling back the formula-theory image. The range formula is
\[
 [\operatorname{Sent}(z)\land\exists p\,G_{e_A}(p,z)]_T.
\]
The graph $G_{e_A}(p,z)\land\operatorname{Sent}(z)$ represents the
sentence-conclusion pullback, so the image--range lemma applies with
the proved partial inverse of its image monomorphism.

\subsection*{(E) Acceptance conditions, non-negotiable}
The integrated proof retains uniform $T$-provable graph equality,
relative totality, finite witness collection, arbitrary slice carriers,
and internal induction for existence and uniqueness. These supply the
appendix's hypotheses. No effective axiomatization or standard soundness
assumption on the ambient $T$ is used.

\subsection*{(F) Not checked / uncertainty log}
The review covers the new realization, the existing definition of $S_T$,
the syntax appendix, and the connected introduction. It does not certify
the remaining proofs of the paper. No Lean files were changed. The
construction asserts initiality in $\BT$; it supplies no claim that
interpretation in every model preserves that initiality. Strongest
confirmed blocker: none. Strongest suspected blocker retained: none.
Whole-paper recommendation: unknown outside the stated review scope.

\section*{Mathematical-writing review}
The governing source was Knuth, Larrabee, and Roberts,
\emph{Mathematical Writing}, Stanford report STAN-CS-89-1193,
Sections 1, 4, and 5. The four-subsection order above was retained.
The comparison below distinguishes source advice from independent
typing conventions of the writing skill.

\begin{longtable}{@{}>{\raggedright\arraybackslash}p{.20\textwidth}>{\raggedright\arraybackslash}p{.40\textwidth}>{\raggedright\arraybackslash}p{.34\textwidth}@{}}
\textbf{Anchor}&\textbf{Issue and implemented correction}&\textbf{Basis}\\\hline
Proposition 5.5&Put coefficient-object and arity hypotheses in the
statement; state the source and target of $K$.&KLR Section 1(5),
and the skill's explicit-typing convention.\\[5pt]
Syntax predicates and objects&Declare $\mathsf{Tm}$ and $\mathsf{Fm}$
as split objects; retain functional notation for the arithmetic
predicates.&KLR Section 1(11), (14); typing convention.\\[5pt]
Finite collection and folds&Distinguish scalar entries, sequence codes,
node records, and evaluation records.&KLR Section 1(14), (17).\\[5pt]
Substitution proof&Define the numerical admissibility predicate;
explain list-occurrence notation; replace set-membership notation
for a categorical variable object.&Typing convention and
definition-order review.\\[5pt]
Theorem 5.9&Use the canonical sentence predicate on formula codes
when intersecting it with the theorem predicate.&Preserves the
common coordinate needed by the formula.\\[5pt]
All new material&Remove duplicate setup, duplicate search proofs,
contributor comments, repeated scope statements, and unnecessary
defensive prose.&The writing skill's economy and manuscript-scope
requirements; KLR Section 1(10), (12), (26).\\
\end{longtable}

The notation audit checked objects, predicates, indices, graph formulas,
slice projections, and range idempotents in reading order. The
introduction now identifies the internal realization and its dependence
on $T$-provable graph equality.

\section*{Sources and document checks}
Robinson and Rosolini, \emph{Categories of partial maps},
Information and Computation 79 (1988), 95--130, pp.~108--109, was
checked in the author-uploaded full text. It is cited for the
functional-$\Sigma_1$ presentation of syntactic partial maps and
provable totality, not as a source for the new realization theorem.
The appendix's existing literature remains in place.

The assembled version 19 manuscript and the separately versioned
appendix were compiled, and their references and layout were checked.
The manuscript has 59 pages; the standalone appendix has 10 pages.
The final source and downloadable files use version 19. The existing
version 18 dependence report describes that earlier version and does
not contain the new Section 5 realization.
\end{document}
